\section{The Guard That Emits Nothing}
\label{sec:ch12-the-guard-that-emits-nothing}

\index{authorization!the second placement|(}%
The obvious correction is to guard the method the entity route calls. The
reference resolver is right there, it is the thing answering, and it takes
one line:

\begin{minted}[highlightlines={3,32}]{csharp}
using HotChocolate.ApolloFederation.Resolvers;
using HotChocolate.ApolloFederation.Types;
using HotChocolate.Authorization;
using HotChocolate.Types.Relay;

namespace Speakers;

[ObjectType<Speaker>]
[Key("id")]
[ReferenceResolver(
    EntityResolver = nameof(ResolveSpeakerReferenceAsync))]
public static partial class SpeakerType
{
    /// <summary>
    /// Loads one speaker by the identifier inside a global node id.
    /// </summary>
    [NodeResolver]
    public static async Task<Speaker?> GetSpeakerAsync(
        int id,
        ISpeakerByIdDataLoader speakerById,
        CancellationToken ct)
        => await speakerById.LoadAsync(id, ct);

    /// <summary>
    /// Answers one representation from the entity route. This is the
    /// method the router reaches every time a client asks a session
    /// for its speaker's name, so it is the busiest resolver in the
    /// graph. The key arrives as the undecoded node id string, and the
    /// type name inside it is checked before the integer is spent.
    /// </summary>
    [GraphQLIgnore]
    [Authorize]
    public static async Task<Speaker?> ResolveSpeakerReferenceAsync(
        string id,
        INodeIdSerializer serializer,
        ISpeakerByIdDataLoader speakerById,
        CancellationToken ct)
    {
        var nodeId = serializer.Parse(id, typeof(int));

        if (nodeId.TypeName != nameof(Speaker))
        {
            return null;
        }

        return await speakerById.LoadAsync(
            (int)nodeId.InternalId!, ct);
    }
}
\end{minted}

\index{authorization!compiles and does nothing}%
It compiles. It compiles with no warning. It guards nothing at all, and the
request from section~\ref{sec:ch12-three-doors-one-column} comes back exactly
as it did before, email and all.

\index{authorization!a caller with a token is told the same thing}%
That is worth separating from a rule that ran and let the caller through.
Present a valid token to that service and the answer is byte for byte the
answer an anonymous caller gets. Nothing is being evaluated. There is no
rule.

\subsection{Look at What Reached the Schema}

\index{authorization!the directive that is not there}%
Export the schema and search it for the directive the attribute is supposed
to write. There are no occurrences. Not on the type, not on a field, and
there is no \code{directive @authorize} definition in the document either.
The attribute produced nothing.

\index{authorization!ApplyPolicy@\code{ApplyPolicy} arrives anyway}%
What did reach the schema is this:

\begin{graphqlsdl}[fontsize=\footnotesize]
"Defines when a policy shall be executed."
enum ApplyPolicy {
  "Before the resolver was executed."
  BEFORE_RESOLVER
  "After the resolver was executed."
  AFTER_RESOLVER
  "The policy is applied in the validation step before the execution."
  VALIDATION
}
\end{graphqlsdl}

\index{authorization!an enum that guards nothing}%
\code{AddAuthorization()} writes that enum whether or not a single field is
guarded, because it is the argument type of a directive the registration
declares. So the only visible trace of authorization in this document is a
type describing when a policy would run, in a schema where no policy runs.
The state reads as configured. It is not.

\subsection{Why the Attribute Had Nowhere to Land}

\index{authorization!field middleware needs a field}%
Hot Chocolate's authorization is field middleware. \code{[Authorize]}
attaches to a field descriptor, and the middleware it installs runs when that
field resolves. A method with no field descriptor gives it nothing to attach
to, and the attribute is dropped without complaint.

\index{GraphQLIgnore@\code{[GraphQLIgnore]}!is what makes the guard inert}%
The reference resolver has no field descriptor, and the attribute directly
above the one I just added is the reason.
Chapter~\ref{ch:first-subgraph} put \code{[GraphQLIgnore]} there so that the
resolver would not be published as a field of \code{Speaker}, which is what
it would otherwise become. That is still right, and it is still doing its
job: there is no \code{resolveSpeakerReference} in the exported schema.

\index{GraphQLIgnore@\code{[GraphQLIgnore]}!the two attributes are in tension}%
So the two attributes on that method want opposite things. One says do not
make this a field. The other only works on fields. They sit on adjacent lines
and neither the compiler nor the schema builder says a word about it.

\index{entities@\code{\_entities}!the resolver is not a field}%
Underneath both of them the reason is the shape of the route. The field being
resolved is \code{\_entities}, declared on \code{Query} and typed
\code{[\_Entity]!}, and \code{\_Entity} is a union of every entity type in
the service. A reference resolver is not a resolver for a field of
\code{Speaker}; it is a thing the entity route calls per representation on
its way to filling that union. There is no field there to hang a rule on
because the only field in sight belongs to the route rather than to the type.

\index{authorization!reported in 2023 and still open}%
This has been in ChilliCream's issue tracker since September 2023, filed by a
user who had exactly this problem and tried exactly these
placements~\autocite{chillicream6546}. It is still open. Michael Staib said
in November of that year that he would look at it, and pointed a month later
at a pull request; that pull request was closed without being merged. I am
not reporting an intention as a fact, and a three-year-old comment is not a
statement about today. What is checkable is what I have just measured: on
16.6.1 this placement still emits nothing.

\index{authorization!the issue never tried the placement that works}%
The useful part of that thread is not the status. It is a comment from
another user, three weeks after the issue opened, working out that no
field-level extension reaches a reference resolver \enquote{since the
resolver itself is not a field}~\autocite{chillicream6546}. That is the
mechanism, and it is right. What the thread never tries is the placement that
does work, which is the subject of section~\ref{sec:ch12-guard-the-data}.
\index{authorization!the second placement|)}%
